% ============================================================================
% CONCLUSION DU CHAPITRE 2 
% source  : partie2/conclusion_partie2.tex
% Version : 2.1 — alignée sur la formalisation Lean4 révisée
% ============================================================================

\section{CONCLUSION DU CHAPITRE~\ref{chap:cadre_formel}}
\label{sec:conclusion_cadre_formel}

\begin{sloppypar}
Le cadre CLO (section~\ref{sec:clo}) a fait de $\Conf$, $\Rel$ et $\Opp$
des grandeurs mesurables. Le
th\'eor\`eme du goulot d'\'etranglement contextuel \'etablit que
$I(V;\Sigma|C) \leq \alpha H(\mathcal{C})$, liant l'opposabilit\'e
accessible \`a la granularit\'e des verdicts et \`a la richesse du contexte
institutionnel. Le cadre distingue la fiabilit\'e de reconstruction
$\Relrec$, contr\^ol\'ee par Fano, de la fiabilit\'e de v\'erification
$\Relver$. Il exige aussi une interpr\'etation admissible pour que $\Opp$ discrimine
les protocoles. Les jeux cryptographiques (IND-CONF,
BIND-REL, EXT-OPP) compl\`etent cette formalisation en \'etablissant le pont
avec les r\'eductions de s\'ecurit\'e standard.

Le Trilemme CRO classique (section~\ref{sec:cro_classique}) a born\'e
ces m\'etriques sous quatre hypoth\`eses explicites sur les limites
institutionnelles, \`a savoir l'extraction born\'ee, le verdict efficace, la r\'egularit\'e
du protocole et le r\'egime de confidentialit\'e $I(M;\Sigma|C) \leq \kappa/\lambda$.
Le produit $\Conf\cdot\Rel\cdot\Opp$ ne peut
d\'epasser $\Delta/(\gamma\lambda) + O(\lambda^{-2})$, tandis que la forme raffin\'ee
identifie le r\'egime ($H(\mathcal{C}) > H^\sharp$) o\`u la capacit\'e
institutionnelle $\Delta$ est la contrainte liante~; l'Hypoth\`ese d'extraction
born\'ee y re\c{c}oit un fondement informationnel exact
(Proposition~\ref{prop:fondement_extraction}). Les syst\`emes d\'eploy\'es ne
sont pas des instances du th\'eor\`eme (Remarque~\ref{rem:portee_cas_etude}).
Ils servent seulement \`a calibrer des capacit\'es implicites,
$\Delta_{\text{Helios}} \approx 56$~bits,
$\Delta_{\text{zk-SNARKs}} \approx 73$~bits, $\Delta_{\text{ECDSA}}
\approx 35$~bits, plausibles au regard de la litt\'erature sur la
rationalit\'e born\'ee.

L'extension post-quantique (section~\ref{sec:cro_pq}) a montr\'e que
cette borne s'\'etend aux adversaires quantiques avec la m\^eme forme,
param\'etr\'ee par la capacit\'e $\Delta_q$. Le coefficient de perte
d'interpr\'etabilit\'e $\eta_q$ n'appara\^it que dans une majoration
affaiblie~; il est nul pour les preuves qui sont des cha\^ines de bits,
y compris les signatures post-quantiques standardis\'ees.

Ces trois acquis, des m\'etriques, une borne classique et son extension
quantique, d\'elimitent ce qui est interdit sans encore indiquer ce qui est
atteignable, question que traite le chapitre~\ref{chap:constructions}.
\end{sloppypar}
